%% Introduction %%

L'aspect monétaire n'a pas seulement un impact sur le traitement documentaire dans le cadre de la commercialisation des archives professionnelles comme nous l'avons démontré plus tôt. A plusieurs reprises le manque de moyen eu un impact sur des politiques documentaires. 

%% Recyclage des bobines %%
\subsection{Économie de matériaux}

A commencer par le tout début de la conservation des documents audiovisuels télévisés. Les moyens manquaient et il était courant de recouper les bobines pour utiliser la bande restante ou pour les faire fondre afin de récupérer l'argent qui couvre les pellicule ou les recycler. C’est une problématique déjà importante pour les archives du cinéma. Pour les archives de la télévision, c’est ce qui motiva Violette Franck à conserver ces documents. C’est cette décision qui donnera naissance à la cinémathèque de l'ORTF et plus tard l’INA. Sans cette héroïne des temps modernes nous pouvons supposer que ces documents auraient connus le même sort que les débuts de la radio en France dont il est périlleux de trouver des traces avant 1944\footcite[14]{jean-michel1986} ou encore des débuts des radios libres, presque impossible à consulter\footcite[16]{jean-michel1986}. De même que dans les débuts du cinéma où les supports de films n'étaient guère considérés que comme des consommables, abandonnés dans des hangars dès l'instant où ils deviennent inutiles, usés jusqu'à la dégradation pour rentabiliser chaque seconde d'image ou découpés, comme le fut la télévision pour recycler les précieuses bandes\footcite[][28]{olmeta2002}. Nous pouvons même nous aventurer aux Etats-Unis où les bandes ont connus cette même menace du coup de ciseau pour économiser les bandes d'enregistrement\footcite[17]{zotero-item-763}. Pour autant beaucoup de documents des débuts de la télévisions ont disparu. C'est ainsi que même aujourd'hui, beaucoup de documents, d'émissions diffusées par l'ORTF manquent toujours. 

%% Stock Gambais %%
\subsection{Le stock Gambais}

Dans ce cadre, le \enquote{stock Gambais} se revêt d'une importance capitale. C'est le nom que l'on donne aux quelques 160 000 boîtes de supports physiques d'archives audiovisuelles retrouvées dans un ancien hangar de l'ORTF à Gambais et aujourd'hui stockés au centre de conservation de Saint-Remy-l'Honoré\footcite[]{theque2024}. Depuis sa découverte au début des années 2000, c'est un travail de longue haleine qui s'est engagé pour dépouiller chaque support et vérifier dans les fonds de l'INA si c'est un document manquant ou si la version retrouvée est de meilleure qualité. De nombreux extraits et programmes complets furent retrouvés ainsi, y compris un fond importants d'émission d'outre-mer, négligés par le traitement documentaire. C'est ainsi que de nouveaux, notices anciennes et récentes se mélangent, entre les notices de documents datant de l'ORTF, documents ayant fait l'objet d'une reprise documentaire et documents issus du stock Gambais, ou repris lors de la découverte d'une meilleure copie. Ainsi, même quand un ou une documentaliste met en garde un chercheur de l'INAthèque contre les documents antérieurs à 1975 pour leur absence de mots clés\footnote{(Voir entretien d'une documentaliste de l'INA sur l'expertise du documentaliste~\ref{ann:transcriptions},\nameref{ann:transcriptions}), p.~\pageref{ann:transcriptions}.}, il y a toujours plusieurs couches d'informations à prendre en compte lors d'une recherche dans les fonds de l'INA. 

%% Coupures budgétaires %%
\subsection{Un budget contraignant}

Pour finir, certaines décisions furent prises pour des raisons financières, notamment dans le cas des types de traitement IND que l'on a cessé d'utiliser pour des raisons de temps et par extension, de budget. Les documentalistes et leur temps sont une ressource précieuse. Il en est finit de recruter en masse des documentalistes comme ce fut le cas à la création du dépôt légal, ce temps doit donc être alloués à ce qui rapporte potentiellement le plus. Ce manque de temps opposé aux quantités de données de plus en plus importantes se ressent dans les politiques documentaires. A partir de 2019, le type de traitement appliqué à un document commence à être déterminé selon son appartenance au fonds dévolu à l'INA. Nous appelons fonds dévolu un fonds incluant les documents pour lesquels la loi a substitué les droits et obligations des sociétés nationales de programme à l'INA (ou le cas échéant à une institution donnée). L'institut a donc tous les droits pour exploiter ces documents. Le fait de déterminer le traitement documentaire sur l'appartenance du document à ces fonds signifie que les documents desquels l'institution peut réaliser un profit aura une description plus appliquée. L'aspect économique prend donc une telle importance qu'il affecte le traitement documentaire du dépôt légal en priorisant une fois de plus l'aspect financier sur le patrimoine. En contrepartie, le traitement de ces émissions repose de plus en plus sur les TGDM, chargés d'ajouter les imports de données, qui voient progressivement s'élargir leur périmètre d'action. Par exemple, récemment, le traitement documentaire de la chaîne France 2 s'est arrêté au profit du traitement des documentalistes de la chaîne télévisées. Il y a certes toujours l'action d'un.e documentaliste et ils utilisent les mêmes référentiels mais les pratiques ne sont pas tout à fait les mêmes, ils ne donnent pas la priorité aux même termes, par dans le même sens car leur mission n'est pas la patrimonialisation en priorité. La nécessité de commercialisation et de rentabilité s'oppose à la communication patrimoniale. 

%% Conclusion %%

L'aspect économique n'influence pas le traitement documentaire de la même façon que le font les usages du traitement documentaire. Il ne contraint les champs descriptifs des documents mais il contraint le périmètre du traitement. Plus les conditions économiques sont favorables, plus le temps des documentalistes est généreusement distribués et plus le périmètres de l'application du traitement documentaire est large et plus il est précis. Cela compromet l'accessibilité aux collections car il n'est pas envisageable de demander à utilisateur des bases de données de l'institution, quel qu'il soit de connaître ces aléas et de comprendre en quoi il affecte son expérience. Cela rejoint le constat que nous avions déjà effectué selon lequel l'accessibilité est compromise par plusieurs facteurs empêchant l'utilisateur de comprendre quelles informations ils peut trouver dans ces bases de données. 